\begin{flushleft}
Účel: chatboty v akademickém prostředí slouží jako 24/7 FAQ, první linie podpory a asistent pro opakující se úkony; pedagogické přínosy při správném návrhu a monitorování.
Doporučený postup: 1) vymezit role a eskalace (citlivé údaje, změny známek => člověk), 2) sesbírat 20–50 scénářů a šablony dialogů, 3) pilot v LMS s fallback, 4) monitorovat logy, upravovat intenty a škálovat governance.
Návrh dialogů: používejte intent‑slot šablony, zobrazujte zdroj a možnost kontaktu; definujte prahy nejistoty pro okamžitou eskalaci.
Technické požadavky: integrace s LMS, autentizace a řízení přístupu, auditní logy, anonymizace a jasné zásady retence.
Monitorování: sledujte přesnost (lidský audit), míru eskalací, spokojenost uživatelů a zavádějte pravidelné lidské kontroly.
Etika a soukromí: transparentně informujte uživatele, nabídněte opt‑out a bezpečný kanál pro citlivé záležitosti.
Kontrolní seznam pro pilot: ověřené FAQ/odpovědi, pravidla eskalace a kontaktní osoby, ochrana dat a mechanismus auditu.
\end{flushleft}